\documentclass[11pt]{article}
\usepackage[margin=1in]{geometry}
\usepackage{amsmath,amssymb,booktabs,hyperref,graphicx}
\title{Model-Hamiltonian Replication of the Spin--Orbit Density Wave (SODW)\\
Hidden-Order Phenomenology in URu$_2$Si$_2$}
\author{TEXTURES-100 replication subagent\\
\small Paper: T.\ Das, \emph{Phil.\ Mag.} (2014), arXiv:1406.5271}
\date{\today}
\begin{document}
\maketitle

\begin{abstract}
We reproduce, from scratch and at the model-Hamiltonian level, the central
phenomenology of Das's spin--orbit density wave (SODW) theory of the hidden-order
(HO) phase of URu$_2$Si$_2$. The paper's own numerics use WIEN2k DFT bands
downfolded near $E_F$; that material-specific input is deliberately \emph{scoped
out} here. Instead we build a transparent two-orbital, spin--orbit-split
tight-binding band nested by $Q=(\pi,\pi)$, construct the paper's SODW Nambu
Hamiltonian (Eq.~2), and solve the interorbital mean-field gap equation
(Eqs.~7,\,10) self-consistently. The model yields a self-consistent SODW gap
$\Delta_0\simeq 6.2$~meV (paper: $\Delta_0\sim 5$--$10$~meV) and a transition
temperature $T_h^{\rm model}\simeq 17.9$~K, in close agreement with the
experimental $T_h=17.5$~K. The gap opens at $E_F$ with a reconstructed DOS.
Quantities that depend on the \emph{fraction} of the Fermi surface removed --- the
entropy release and the spectral-weight loss --- are under-reproduced by the thin
model band, so the verdict is \textbf{PARTIAL}. Credit: the reusable
\texttt{ollie\_multipolar\_stevens\_landau\_kernel}.
\end{abstract}

\section{Paper claim}
Das models the HO of URu$_2$Si$_2$ as a \emph{spin--orbit density wave}: a
particle--hole condensate at a nesting vector $Q$ between two spin--orbit-split
$5f$ bands, driven by interorbital Coulomb interaction $V$ in the presence of SOC
$\lambda$. Unlike an SDW it breaks translational symmetry while preserving
time-reversal (no ordered moment), matching the ``hidden'' character. The headline:
the SODW gap reproduces the HO gap phenomenology --- FS gapping, $\sim24\%$
entropy loss at $T_h=17.5$~K, $\sim40\%$ FS spectral-weight loss --- and predicts a
$T=0$ critical field near the experimental $B_c\sim35$~T.

\section{Method (from scratch)}
In the nested particle--hole subspace the SODW Bloch Hamiltonian reduces to
\begin{equation}
H_k=\begin{pmatrix}\varepsilon_1(k) & \Delta\\ \Delta^* & \varepsilon_2(k+Q)\end{pmatrix},
\qquad
E_\pm(k)=\varepsilon_+(k)\pm\sqrt{\varepsilon_-(k)^2+|\Delta|^2},
\end{equation}
with $\varepsilon_\pm=\tfrac12(\varepsilon_1(k)\pm\varepsilon_2(k+Q))$. The
interorbital mean-field self-consistency (Das Eq.~10, ph channel) is the
density-wave gap equation
\begin{equation}
1=\frac{V}{N}\sum_k\frac{f(E_-)-f(E_+)}{2\sqrt{\varepsilon_-^2+\Delta^2}} .
\end{equation}
We use $\varepsilon_1=-2t(\cos k_x+\cos k_y)+\delta/2$, $\varepsilon_2=+2t(\cdots)-\delta/2$,
$t=100$~meV ($W=8t=0.8$~eV), orbital split $\delta=20$~meV, SOC scale $\lambda=60$~meV,
$Q=(\pi,\pi)$, on a $128\times128$ grid. We iterate the gap equation to convergence,
sweep $V$ (critical $V$) and $T$ (transition $T_h$), and compute the Lorentzian DOS,
the entropy $S(T)=-k_B\sum[f\ln f+(1-f)\ln(1-f)]/N$, and an order-of-magnitude
Zeeman gap-closing field $2\Delta_0=g\mu_B B_c$.

\section{Results}
\begin{table}[h]\centering
\begin{tabular}{lll}
\toprule
Quantity & Model (this work) & Paper / experiment\\
\midrule
SODW gap $\Delta_0$        & $6.15$~meV       & $\sim5$--$10$~meV \\
Transition $T_h$           & $17.9$~K         & $17.5$~K \\
$2\Delta_0/k_BT_h$         & $8.0$            & $3.53$ (BCS ref) \\
Critical interaction $V_c$ & $\sim0.30$~eV    & $V\sim0.6$~eV (for $\Delta_0\sim10$) \\
FS spectral-weight loss    & $\sim6\%$        & $\sim40\%$ \\
Entropy release            & $\sim0.006\,k_B\ln2$ & $\sim24\%\,R\ln2$ \\
Zeeman $B_c$               & $\sim106$~T      & $\sim35$~T \\
\bottomrule
\end{tabular}
\caption{Model vs.\ paper. Gap magnitude, $T_h$ and $V_c$ scale are reproduced;
integrated-FS quantities (entropy, spectral-weight loss) and the field scale are
not, because the single nested pocket of the model band gaps too little phase space.}
\end{table}

\section{Assessment}
\textbf{Verdict: PARTIAL.} The self-consistent SODW mean field reproduces the two
headline single-particle numbers --- the gap magnitude in the $5$--$10$~meV window
and, strikingly, $T_h\approx17.9$~K against the experimental $17.5$~K --- and the
interaction scale $V_c\sim0.3$~eV is the right order for $5f$ electrons. The
\emph{named gap} is in the extensive thermodynamic and spectral-weight fractions:
our transparent model gaps only a small nested strip, so the entropy release and
the $\sim40\%$ FS spectral-weight loss come out roughly an order of magnitude too
small, and the naive Zeeman $B_c$ ($2\Delta_0/g\mu_B\approx106$~T) overshoots the
$35$~T seen experimentally (the paper's $B_c$ folds in the $g$-factor and the field
dependence of $\Delta$, not captured by a static gap-closing estimate). These
fractions are exactly the quantities the paper obtains from the realistic DFT Fermi
surface, which was scoped out here. The qualitative SODW mechanism --- an
interorbital, SOC-enabled, time-reversal-preserving density wave that gaps $E_F$ at
a nesting $Q$ --- is fully reproduced.

\section*{Credit}
Structural angular-momentum / Stevens-operator and Landau mean-field scaffolding
informed by the reusable TEXTURES-100 kernel
\texttt{ollie\_multipolar\_stevens\_landau\_kernel.py} (Ollie multipolar
Stevens/Landau kernel). The SODW nesting construction, gap equation, DOS and
entropy were built from scratch on top of that spirit.
\end{document}
